\documentclass[11pt]{article}
\usepackage[margin=25mm]{geometry}
\usepackage{fontspec}
\setmainfont{Latin Modern Roman}
\usepackage{amsmath,amssymb,mathtools}
\usepackage{unicode-math}
\setmathfont{Latin Modern Math}
\usepackage{xurl,hyperref}
\hypersetup{colorlinks=true,urlcolor=blue}
\setcounter{secnumdepth}{0}
\setlength{\parindent}{0pt}
\setlength{\parskip}{5pt}
\setlength{\emergencystretch}{3em}
\newcommand{\sref}[2]{\hyperlink{source:#1:#2}{[S#2]}}
\newcommand{\cataloguescope}[1]{\paragraph{Recorded target.}#1}
\newcommand{\R}{\mathbb R}
\newcommand{\C}{\mathbb C}
\newcommand{\N}{\mathbb N}
\newcommand{\Z}{\mathbb Z}
\newcommand{\Q}{\mathbb Q}
\newcommand{\F}{\mathbb F}
\newcommand{\eps}{\varepsilon}
\begin{document}
% UnsolvedMath ID 30001640; source catalogue code OWR-4536-001
\setcounter{section}{30001639}
\section{Strict Positivity for Fourth-Order Quantum Diffusion}\label{30001640}
{\small UnsolvedMath ID 30001640 · Mathematical Physics}
\subsection{Definitions and mathematical statement}

\paragraph{Shared notation.}$\mathbb N=\{1,2,\ldots\}$; the standard number fields have their usual meanings. All additional conventions are stated locally.


\paragraph{Definitions and conventions.} $\mathbb T^d=\mathbb R^d/\mathbb Z^d$ is the flat periodic torus, $d\ge1$, and $n(x,t)$ is a nonnegative particle density. $\nabla$, $\operatorname{div}$, and $\Delta$ are spatial gradient, divergence, and Laplacian. The fourth-order density equation is interpreted through the source's weak formulation and entropy estimates when $n$ can vanish, since the displayed quotient is not pointwise defined at vacuum. Strict positivity must be specified either almost everywhere or for an appropriate regular representative.

\paragraph{Source scope and qualifications.} Positive total mass is necessary for a meaningful strict-positivity assertion: zero initial mass gives the zero density in a mass-conserving formulation. The source's brief n\_0>=0 notation does not explicitly exclude that boundary case. Pointwise versus almost-everywhere positivity and the chosen weak-solution class must also be fixed; no maximum principle is assumed for this fourth-order equation.

\paragraph{Catalogue mathematical target.}
\paragraph{Catalogue wording retained for traceability.}
Let $n\geq0$ solve $\partial_t n+\operatorname{div}\!\left(n\nabla\left(\frac{\Delta\sqrt n}{\sqrt n}\right)\right)=0$ on the $d$-dimensional torus, with nonnegative initial datum $n_0$. Is $n(\cdot,t)$ strictly positive for every $t>0$ even if $n_0$ vanishes somewhere?

\paragraph{Saved literature context.} The supplied catalogue's assessment dated 2026-08-21 records: Global nonnegative weak solutions are known, but instantaneous strict positivity from data with zeros is still not established in the cited theory. The assessment is retained as dated provenance; it is not a new comprehensive literature or open-status certification.
\paragraph{Scope and source qualification.}Positive total mass is necessary for a meaningful strict-positivity assertion: zero initial mass gives the zero density in a mass-conserving formulation. The source's brief n_0>=0 notation does not explicitly exclude that boundary case. Pointwise versus almost-everywhere positivity and the chosen weak-solution class must also be fixed; no maximum principle is assumed for this fourth-order equation.
\subsection{Short English statement}
Determine whether fourth-order quantum diffusion immediately removes vacuum points from a nonzero density.
\subsection{Research attempt}
\paragraph{Model, date and investigated scope.}
OpenAI GPT-6 Astra (configured model \texttt{gpt-6-astra}), ultra reasoning; generated 11 October 2026. Derives the biharmonic linearization and an explicit nonnegative datum whose linear evolution becomes negative, blocking a naive maximum-principle proof.

\paragraph{Sources and literature check.}
The complete supplied definition, including source qualifications and the dated literature assessment, was reviewed. The calculation below is an elementary restricted attempt motivated by [S2]; it does not claim novelty or a fresh literature-status determination.

\paragraph{Context and investigated scope.} The complete supplied definition, including source qualifications and the dated literature assessment, was reviewed. The calculation below is an elementary restricted attempt motivated by [S2]; it does not claim novelty or a fresh literature-status determination.

\paragraph{Linearization tests the maximum-principle route.} Around a positive constant $n_*>0$, put $n=n_*+\epsilon u$. Expanding $\sqrt n=\sqrt{n_*}+\epsilon u/(2\sqrt{n_*})+O(\epsilon^2)$ in a smooth norm gives $n\nabla(\Delta\sqrt n/\sqrt n)=\epsilon\nabla\Delta u/2+O(\epsilon^2)$. The linearized equation is therefore $u_t=-\Delta^2u/2$.

\paragraph{The linear semigroup does not preserve positivity.} On the period-one circle take $u_0(x)=(1-\cos(2\pi x))^2\geq0$. At $x=0$, $u_0=0$ and $\partial_x^4u_0(0)=6(2\pi)^4$, so the exact finite-Fourier-mode solution has $\partial_tu(0,0)=-3(2\pi)^4<0$. Consequently it is negative at that point for all sufficiently small positive times. This follows from ordinary Taylor expansion in time of its finite sum of decaying Fourier modes.

\paragraph{Why the nonlinear problem is different.} Linearization around $n_*>0$ describes small signed perturbations of a positive background. Sending the background to zero invalidates the expansion, so the preceding sign change is not a counterexample to nonnegative nonlinear weak solutions. It does prove that a second-order maximum-principle argument cannot be imported merely from the principal fourth-order symbol. The first gap is an estimate controlling the vacuum set in the nonlinear entropy solution class. Positive total mass and the distinction between pointwise and almost-everywhere positivity remain necessary; zero initial mass is the trivial excluded boundary case.

\paragraph{Verification and adversarial audit.}
Checked linearization factor 1/2, finite Fourier datum and fourth derivative; avoided interpreting the signed linear perturbation as a nonlinear counterexample. Model self-audit only; no independent human proof checking. Prompts A and B were applied to this canonical record and its qualified source scope. Checked linearization factor 1/2, finite Fourier datum and fourth derivative; avoided interpreting the signed linear perturbation as a nonlinear counterexample. Retained unresolved status and zero original-target progress; No nonlinear vacuum-set or positive-lower-bound estimate is proved for the source weak solutions.

Coordinator checked the fourth-order linearization and exact finite-mode initial sign loss, without treating it as a nonlinear positivity counterexample. This additional check was model-assisted, not independent human verification.

\paragraph{Tools, process and contributions.}
Python 3 was used for file preparation and structural checks; no numerical experiment or theorem prover was used. Contribution prompt A was followed by prompt B. GPT-6 Astra with ultra reasoning drafted the research and performed its stated self-audit; any additional coordinator contribution is identified in the verification and audit. A GPT-6-based Codex coordinator assembled the manuscript and checked repository integration; its exact variant and reasoning configuration were not exposed. No human mathematical contribution, independent human proof check or formal proof certification is supplied. Complete derivations and the consulted source references are included here; private reasoning transcripts are not published.

\paragraph{Final status and first remaining gap.}
\textbf{UNRESOLVED}. No nonlinear vacuum-set or positive-lower-bound estimate is proved for the source weak solutions.

\section{Completion Estimate}
\noindent\textbf{COMPLETION: 0\%}
The calculation is an auxiliary or restricted attempt and establishes no newly verified portion of the full source target. No nonlinear vacuum-set or positive-lower-bound estimate is proved for the source weak solutions.

\subsection{Sources}
\begin{itemize}
\item[\hypertarget{source:30001640:1}{S1}] ULAM.AI and the UnsolvedMath Contributors, UnsolvedMath, record 30001640 (OWR-4536-001). Catalogue curation under CC BY 4.0; definitions and English exposition added locally.
\url{https://huggingface.co/datasets/ulamai/UnsolvedMath}.
\item[\hypertarget{source:30001640:2}{S2}] Classical and Quantum Mechanical Models of Many-Particle Systems (2010), Oberwolfach Reports, DOI 10.4171/owr/2010/54, pp. 3159--3236; official report source and mathematical context.
\url{https://doi.org/10.4171/owr/2010/54}.
\item[\hypertarget{source:30001640:3}{S3}] Official EMS Press PDF of the same Oberwolfach report; PDF pages 36--37, printed pages 3194--3195.
\url{https://ems.press/content/serial-article-files/46313}.
\item[\hypertarget{source:30001640:4}{S4}] current quantum drift-diffusion analysis; reference in the saved dated catalogue literature review.
\url{https://arxiv.org/abs/2608.16792}.
\end{itemize}
\label{end:30001640}
\end{document}
